% Generated draft: requires mathematical and source review.
\documentclass[11pt]{article}
\usepackage[margin=25mm]{geometry}
\usepackage{fontspec}
\setmainfont{Latin Modern Roman}
\usepackage{amsmath,amssymb,mathtools,unicode-math}
\setmathfont{Latin Modern Math}
\usepackage{enumitem,needspace,xurl,hyperref}
\setlength{\parindent}{0pt}
\setlength{\parskip}{5pt}
\setlength{\emergencystretch}{3em}
\newcommand{\N}{\mathbb{N}}
\newcommand{\Z}{\mathbb{Z}}
\newcommand{\Q}{\mathbb{Q}}
\newcommand{\R}{\mathbb{R}}
\newcommand{\C}{\mathbb{C}}
\newcommand{\F}{\mathbb{F}}
\newcommand{\A}{\mathbb{A}}
\newcommand{\PP}{\mathbb{P}}
\newcommand{\E}{\mathbb{E}}
\newcommand{\eps}{\varepsilon}
\newcommand{\dd}{\,\mathrm{d}}
\begin{document}
\section*{UnsolvedMath ID: 700891}
\textbf{Source problem code:} catalogue entry
\subsection{Definitions and mathematical statement}
A formation of finite groups is a possibly empty class of finite groups, closed under isomorphisms, homomorphic images, and subdirect products of finitely many groups with at least one factor. A subgroup of a direct product is subdirect if every coordinate projection is surjective. A finite group is soluble if its derived series terminates at the trivial group. A soluble formation is a formation all of whose members are soluble.

A formation \(\mathfrak A\) is a maximal proper subformation of \(\mathfrak B\) if \(\mathfrak A\subsetneq\mathfrak B\) and there is no formation \(\mathfrak K\) satisfying \(\mathfrak A\subsetneq\mathfrak K\subsetneq\mathfrak B\). Following the supplied definition, a formation \(\mathfrak H\) has length \(t\) if there exists a chain
\[
\varnothing=\mathfrak H_0\subsetneq\mathfrak H_1\subsetneq\cdots\subsetneq\mathfrak H_t=\mathfrak H
\]
in which \(\mathfrak H_{i-1}\) is a maximal proper subformation of \(\mathfrak H_i\) for every \(1\leq i\leq t\). This is an existence condition on a chain; no uniqueness of its length is assumed. Adopt the zero-step convention for \(\mathfrak H=\varnothing\), a boundary case not explicitly discussed in the source.

Let \(\mathcal L_{\leq4}\) be the collection of soluble formations admitting such a chain for some integer \(0\leq t\leq4\), ordered by inclusion. The source asks whether the lattice of these formations is distributive. Interpreted as a question about this ordered collection, the question is whether it is a lattice whose meet and join satisfy, for all \(\mathfrak A,\mathfrak B,\mathfrak C\in\mathcal L_{\leq4}\),
\[
\begin{aligned}
\mathfrak A\wedge(\mathfrak B\vee\mathfrak C)
&=(\mathfrak A\wedge\mathfrak B)\vee(\mathfrak A\wedge\mathfrak C),\\
\mathfrak A\vee(\mathfrak B\wedge\mathfrak C)
&=(\mathfrak A\vee\mathfrak B)\wedge(\mathfrak A\vee\mathfrak C).
\end{aligned}
\]
Here meet and join mean greatest lower bound and least upper bound in the stated ordered collection. The source calls this collection a lattice without explaining its operations or their existence. In the ambient lattice of all soluble formations, meet is intersection and join is the least formation containing both operands; whether those operations are intended to restrict to \(\mathcal L_{\leq4}\) requires clarification.
\subsection{Short English statement}
A. N. Skiba asks whether soluble formations obtainable from the empty formation by at most four successive maximal-subformation inclusions form a distributive lattice. The source leaves the lattice operations implicit and defines length through the existence of a chain.
\subsection{Sources}
\begin{itemize}
\item Supplied local UnsolvedMath registry record: permanent ID \(700891\), exact source problem code \texttt{catalogue entry}, title \emph{Kourovka Notebook Problem 12.73}. Upstream revision: \texttt{d371ee917ac9647fbb06cd6fc529deec6f3dda82}. The record labels the problem \emph{open}; this is supplied source metadata, not an independently verified status assessment.
\item Supplied bibliographic reference: \emph{The Kourovka Notebook (Unsolved Problems in Group Theory)}, original Problem 12.73 in the 12th issue (1992), as recorded in the 21st issue edition (2026), \texttt{arXiv:1401.0300v48}, \url{https://arxiv.org/abs/1401.0300v48}. The supplied collection description names E. I. Khukhro and V. D. Mazurov as editors. The cited edition has not been independently consulted here.
\item Attribution supplied in the record: proposed by A. N. Skiba in 1992.
\end{itemize}
\end{document}
